\documentclass{hw-template}
\usepackage{cancel}

\title{MA 587 Homework 2}
\author{Kyle Hansen}
\date{\today}

\begin{document}

\newcommand{\hn}[1]{\left\| {#1} \right\|_{H^k}}\
\newcommand{\iph}[2]{\left({#1}, {#2}\right)_{H^k}}
\newcommand{\ipl}[2]{\left\langle{#1}, {#2}\right\rangle_{L_2(\Omega)}}
\newcommand{\da}{D^{\alpha}}

\maketitle

\section{Error bound}

\subsection{} Consider a single cell of the problem domain, $L_j = [x_{j-1}, x_j]$, and choose some $z \in (x_{j-1}, x_j)$. The quadratic polynomial that interpolates the points $\{u(x_{j-1}), u(z), u(x_{j})\}$ is given by

\begin{equation} \label{eq:quadratic}
    \hat{u}_z(x) = u_{j-1} + \frac{u_j - u_{j-1}}{h}x + A_z(x-x_j)(x-x_{j-1}),
\end{equation}

where $A_z$ is some constant dependent on the choice of $z$, $u_j = u(x_j)$, and $h = x_j - x_{j-1}$. This can also be rewritten as 

\begin{equation}
    \hat{u}_z(x) = \tilde{u}(x) + A_z(x-x_j)(x-x_{j-1}),
\end{equation}

where $\tilde{u}$ is the linear interpolation of $u$. Then the error $g(x) = \hat{u}(x) - u(x)$ is zero in at least three points: $g(x_{j-1}) = g(z) = g(x_j) = 0$. By Rolle's theorem, since $g$ has three zeros, there are two points $\xi_{1, z} \in (x_{j-1}, z)$ and $\xi_{2, z} \in (z, x_j)$, for which $g^\prime(\xi_{1, z}) = g^\prime(\xi_{2, z}) = 0$. Applying Rolle's theorem once more, there is another point $\xi_{z} \in (\xi_{1, z}, \xi_{2, z})$ such that $g^{\prime\prime}(\xi_z) = 0$. That is, 

\begin{equation}
    g^{\prime\prime}(\xi_z) = \hat{u}^{\prime\prime}(\xi_z) - u^{\prime\prime}(\xi_z), 
\end{equation}

so, using \autoref{eq:quadratic}, 

\begin{gather}
    u^{\prime\prime}(\xi_z) = 2A_z \\
    \frac{u^{\prime\prime}(\xi_z)}{2} = A_z.
\end{gather}

This value is bounded by

\begin{equation}\label{eq:bound}
    \left|\frac{u^{\prime\prime}(\xi_z)}{2} \right| \leq \max_{a \in [x_{j-1}, x_j]}  \left|\frac{u^{\prime\prime}(a)}{2}\right| \leq \max_{y \in [0, 1]}  \left|\frac{u^{\prime\prime}(y)}{2}\right|
\end{equation}

Then, for any $x \in [x_{j-1}, x_j]$, 

\begin{align}
    u(x) &= \tilde{u}(x) + A_x(x-x_j)(x-x_{j-1})\\
    u(x) &= \tilde{u}(x) + \frac{u^{\prime\prime}(\xi_x)}{2}(x-x_j)(x-x_{j-1})\\
    u(x) - \tilde{u}(x) &= \frac{u^{\prime\prime}(\xi_x)}{2}(x-x_j)(x-x_{j-1}),
\end{align}

so

\begin{equation}
    \left|u(x) - \tilde{u}(x)\right| = \left|\frac{u^{\prime\prime}(\xi_x)}{2}(x-x_j)(x-x_{j-1}) \right|.
\end{equation}

By submultaplicativity, 

\begin{equation}
    \left|u(x) - \tilde{u}(x)\right| \leq \left|\frac{u^{\prime\prime}(\xi_x)}{2}\right| \cdot \left|(x-x_j)(x-x_{j-1}) \right|.
\end{equation}

Applying the triangle inequality to the $(x-x_j)(x-x_{j-1})$ term and applying \autoref{eq:bound} to the $\frac{u^{\prime\prime}}{2}$ term, this is bounded by

\begin{equation}
   \left|u(x) - \tilde{u}(x)\right| \leq  \max_{y \in [0, 1]}  \left|\frac{u^{\prime\prime}(y)}{2}\right| \cdot \left( \left| x-x_j\right| + \left|x-x_{j-1} \right| \right),
\end{equation}

or 

\begin{gather}
    \left|u(x) - \tilde{u}(x)\right| \leq  \max_{y \in [0, 1]}  \left|\frac{u^{\prime\prime}(y)}{2}\right| \cdot h\\
    \left|u(x) - \tilde{u}(x)\right| \leq  \max_{y \in [0, 1]}  \left|{u^{\prime\prime}(y)}\right| \cdot h.
\end{gather}

\subsection{} The inequality given is

\begin{equation}
    \norm{(u - u_h)^\prime} \leq h \cdot \max_{y \in [0, 1]}  \left|{u^{\prime\prime}(y)}\right|.
\end{equation}


The Cauchy-Schwarz Inequality states that $|\ip{\vec{x}}{\vec{y}}| \leq \norm{\vec{x}} \cdot \norm{\vec{y}}$, where $\ip{\vec{x}}{\vec{y}} = \int_0^1 \vec{x}\vec{y}dx$ is the inner product, and the norm $\norm{\cdot}$ is $\ip{\cdot}{\cdot}^{1/2}$ Let $\vec{x} = (u - u_h)^\prime$ and let $\vec{y} = 1$:

\begin{equation}
    |\ip{(u - u_h)^\prime}{1}| \leq \norm{(u - u_h)^\prime} \cdot \norm{1} \leq h \cdot \max_{y \in [0, 1]}  \left|{u^{\prime\prime}(y)}\right|,
\end{equation}

and, expanding the inner product definition,

\begin{equation}
    \left|\int_0^1 (u - u_h)^\prime dx \right| \leq h \cdot \max_{y \in [0, 1]}  \left|{u^{\prime\prime}(y)}\right|.
\end{equation}

Using a change of variables,

\begin{equation}
    \left|\int_0^x (u - u_h)^\prime (s) ds \right| \leq \left| \int_0^1 (u - u_h)^\prime (s) ds \right| \leq h \cdot \max_{y \in [0, 1]}  \left|{u^{\prime\prime}(y)}\right|
\end{equation}

for $x \in [0, 1]$, which becomes

\begin{equation}
     \left|\left[ u(x) - u_h(x) \right] - \cancel{\left[ u(0) - u_h(0) \right]} \right| \leq h \cdot \max_{y \in [0, 1]}  \left|{u^{\prime\prime}(y)}\right|
\end{equation}

\begin{equation}
     \left| u(x) - u_h(x) \right| \leq h \cdot \max_{y \in [0, 1]}  \left|{u^{\prime\prime}(y)}\right|.
\end{equation}

\section{Formulations}

\subsection{} First, to show that (V)$\implies$ (M):

Let $u^v$ be a solution to (V), so

\begin{equation}\label{eq:variational}
    a(u^v, v) = L(v) \forall v \in \mathbb{V}.
\end{equation}

Then let $w \in \mathbb{V}$, and consider

\begin{equation}
    F\left( u^v + w \right) = \frac{1}{2}a\left(u^v + w, u^v + w \right) - L\left( u^v + w\right).
\end{equation}

Since $a$ is bilinear and $L$ is linear, this can be rewritten

\begin{align}
    F\left( u^v + w \right) &= \frac{1}{2}  a(u^v, u^v) + a(u^v, w) + \frac{1}{2}a(w, w) - L\left( u^v \right) - L(w)\\
    F\left( u^v + w \right) &= \frac{1}{2}  a(u^v, u^v)+ \frac{1}{2}a(w, w)  + \left[ a(u^v, w) - L(w) \right] - L\left( u^v \right).
\end{align}

From \autoref{eq:variational}, the bracketed term is zero, and 

\begin{align}
    F\left( u^v + w \right) &= \frac{1}{2}  a(u^v, u^v)+ \frac{1}{2}a(w, w)   - L\left( u^v \right)\\
    F\left( u^v + w \right) &= F(u^v) + \frac{1}{2}a(w, w) \geq F(u^v).
\end{align}

Since the bilinear $a(w,w)$ is strictly non-negative, this inequality applies. For any $w$, $F(u^v)\leq F(u^v + w)$, so for any $v$, $F(u^v)\leq F(v)$ and (M) is satistified, and thus (V) $\implies$ (M).

Next, to show that (M) $\implies$ (V), consider $u^m$ that satisfies (M). Then $u^m$ minimizes $F()$. Consider the scalar function

\begin{equation}
    f(\epsilon) = F(u^m + \epsilon v)
\end{equation}

for some $v$. Since $u_m$ minimizes $F$, then $\epsilon = 0$ minimizes $f$, so

\begin{equation}
    \left.\frac{\partial f}{\partial \epsilon} \right|_{\epsilon = 0} = 0.
\end{equation}

Expanding this yields

\begin{gather}
     \left.\frac{\partial }{\partial \epsilon} \left( F(u^m + \epsilon v) \right) \right|_{\epsilon = 0} =0 \\
     \left.\frac{\partial }{\partial \epsilon}\left( \frac{1}{2} a(u^m + \epsilon v, u^m + \epsilon v) - L(u^m + \epsilon v) \right) \right|_{\epsilon = 0} =0\\
     \left.\frac{\partial }{\partial \epsilon}\left( \frac{1}{2} a(u^m, u^m) + a(u^m, \epsilon v) + \frac{1}{2}a(\epsilon v, \epsilon v) - L(u^m) - L(\epsilon v) \right) \right|_{\epsilon = 0}=0\\
     \left.\frac{\partial }{\partial \epsilon}\left( \frac{1}{2} a(u^m, u^m) + \epsilon a(u^m,  v) + \frac{\epsilon^2}{2}a( v,  v) - L(u^m) - \epsilon L( v) \right) \right|_{\epsilon = 0}=0\\
     \left. \vphantom{\frac{}{}} a(u^m, v) + \epsilon a(v, v) - L(v) \right|_{\epsilon = 0}=0\\
      a(u^m, v)  - L(v) =0
\end{gather}

since this applies for any $v \in \mathbb{V}$, (V) is satisfied and thus (M) $\implies$ (V). Since (M) $\implies$ (V) and (V) $\implies$ (M), (V)$\iff$ (M).

\subsection{} Since $\mathbb{V}_h$ is a Hilbert space and $a$ is symmetric, linear, and coercive, and from part a, {(V) $\iff$ (M)}, so the variational solution $u_h$ also satisfies

\begin{align}
    F(u_h) &\leq F(v) \quad \forall\; v \in \mathbb{V}_h\\
    2F(u_h) &\leq 2F(v) \\
    a(u, u) + 2F(u_h) &\leq a(u, u) +  2F(v) \\
    a(u, u) + a(u_h, u_h) - 2L(u_h) &\leq a(u, u) +  a(v, v) - 2L(v).
\end{align}

Since $u$ satisfies (V), $a(u, w) = L(w)$ for any $w$ in $\mathbb{V}$, including members of the subset $\mathbb{V}_h$, so 

\begin{equation}
    a(u, u) + a(u_h, u_h) - 2a(u, u_h) \leq a(u, u) +  a(v, v) - 2a(u, v)\quad \forall\; v \in \mathbb{V}_h.
\end{equation}

This can be rewritten, using the bilinearity of $a$, as

\begin{align}
    a(u-u_h, u-u_h) &\leq a(u-v, u-v)\\
    \norm{u-u_h}_a^2 &\leq \norm{u-v}_a^2\\
    \norm{u-u_h}_a &\leq \norm{u-v}_a \quad \forall\; v \in \mathbb{V}_h.
\end{align}

In other words, $u_h$ is the best approximation of $u$ of any function in the discretized set $\mathbb{V}_h$, using $a$-norm error.

\section{Basis Functions}

\subsection{} The three points form a triangle as long as they are not collinear, or if 

\begin{equation}
    [N_1 - N_2]^T[N_1 - N_3] \neq |N_1 - N_2| \cdot |N_1 - N_3|.
\end{equation}

\subsection{}

Denoting the basis functions using

\begin{equation}
    \phi_j(x, y) = a_j x + b_j y + c_j,
\end{equation}

the condition can be notated as a linear system

\begin{equation}
    \begin{bmatrix}
        1 & 0 & 0\\
        1 & h & 0\\
        1 & 0 & h
    \end{bmatrix} \begin{bmatrix}
        c_1 & c_2 & c_3\\
        a_1 & a_2 & a_3\\
        b_1 & b_2 & b_3
    \end{bmatrix} = \begin{bmatrix}
        1 & 0 & 0\\
        0 & 1 & 0\\
        0 & 0 & 1
    \end{bmatrix}
\end{equation}

Then the coefficients are simply the entries of 

\begin{equation}
    A = \begin{bmatrix}
        1 & 0 & 0\\
        1 & h & 0\\
        1 & 0 & h
    \end{bmatrix}.
\end{equation}

This matrix $A^{-1}$ is easily found, since $A$ is lower triangular, to be

\begin{equation}
    A^{-1} = \begin{bmatrix}
        1 & 0 & 0\\
        \frac{-1}{h} & \frac{1}{h} & 0\\
        \frac{-1}{h} & 0 & \frac{1}{h}
    \end{bmatrix}.
\end{equation}

Then the basis functions are 

\begin{align}
    \phi_1(x, y) &= 1 - \frac{x}{h} + \frac{y}{h}\\
    \phi_2(x, y) &= \frac{x}{h}\\
    \phi_3(x, y) &= \frac{-1}{h}+ \frac{y}{h}.
\end{align}

Since a 2d linear function can be represented using coefficients $a, b, c$, any basis that spans $\mathbb{R}^3$ can be a basis for 2d linear functions. The columns of $A^{-1}$ are mutually orthogonal, and do span $\mathbb{R}^3$, so these basis functions $\phi_j$ do form a basis for 2d linear functions.

\subsection{} First, the gradients of the basis functions are 

\begin{align}
    \nabla\phi_1(x, y) &= \begin{bmatrix}
        \frac{-1}{h} & \frac{1}{h}
    \end{bmatrix}^T\\
    \nabla\phi_2(x, y) &= \begin{bmatrix}
        \frac{1}{h} & 0
    \end{bmatrix}^T\\
    \nabla\phi_3(x, y) &= \begin{bmatrix}
        0 & \frac{1}{h}
    \end{bmatrix}^T.
\end{align}

Then, the stiffness matrix entries are 

\begin{align}
    a(\phi_1, \phi_1) &= \int_K \nabla \phi_1 \cdot \nabla\phi_1 dx\\
                      &= \int_K  \left(\frac{1}{h} \right)^2 + \left(\frac{1}{h} \right)^2dx\\
                      &= \frac{2}{h^2} \int_K dx\\
                      &=1
\end{align}

\begin{align}
    a(\phi_1, \phi_2) &= \int_K \nabla \phi_1 \cdot \nabla\phi_2 dx\\
                      &= \int_K  \left(\frac{1}{h} \cdot \frac{-1}{h} \right) dx\\
                      &= \frac{-1}{h^2} \int_K dx\\
                      &= \frac{-1}{2}
\end{align}

\begin{align}
    a(\phi_1, \phi_3) &= \int_K \nabla \phi_1 \cdot \nabla\phi_3 dx\\
                      &= \int_K  \left(\frac{1}{h} \cdot \frac{-1}{h} \right) dx\\
                      &= \frac{-1}{h^2} \int_K dx\\
                      &= \frac{-1}{2}
\end{align}

\begin{align}
    a(\phi_2, \phi_2) &= \int_K \nabla \phi_2 \cdot \nabla\phi_2 dx\\
                      &= \int_K  \left(\frac{1}{h} \cdot \frac{1}{h} \right) dx\\
                      &= \frac{1}{h^2} \int_K dx\\
                      &= \frac{1}{2}
\end{align}

\begin{align}
    a(\phi_2, \phi_3) &= \int_K \nabla \phi_2 \cdot \nabla\phi_3 dx\\
                      &= \int_K  0 dx\\
                      &= 0
\end{align}

\begin{align}
    a(\phi_3, \phi_3) &= \int_K \nabla \phi_3 \cdot \nabla\phi_3 dx\\
                      &= \int_K  \left(\frac{1}{h} \cdot \frac{1}{h} \right) dx\\
                      &= \frac{1}{h^2} \int_K dx\\
                      &= \frac{1}{2},
\end{align}

The stiffness matrix is symmetric since $a$ is bilinear, so the stiffness matrix $A$ matches the one given in the problem description.

\section{Sobolev Space}

\subsection{}

\begin{align}
    \hn{u} &= \sqrt{\iph{u}{u}}\\
    \hn{u}^2 &= \iph{u}{u}\\
             &= \sum_{|\alpha|\leq k} \ipl{D^{\alpha}u}{D^{\alpha}u}\\
             &= \sum_{|\alpha|\leq s} \ipl{D^{\alpha}u}{D^{\alpha}u} + \sum_{s < |\alpha|\leq k} \ipl{D^{\alpha}u}{D^{\alpha}u}\\
             &= \norm{u}_{H^s}^2 + \sum_{s < |\alpha|\leq k} \ipl{D^{\alpha}u}{D^{\alpha}u}.\label{eq:last}
\end{align}

Since the $L_2$ inner product is strictly non-negative, the rightmost term in \autoref{eq:last} is nonnegative, so

\begin{align}
    \norm{u}_{H^s}^2 + \sum_{s < |\alpha|\leq k} \ipl{D^{\alpha}u}{D^{\alpha}u} = \hn{u}^2\\
    \sqrt{\norm{u}_{H^s}^2 + \sum_{s < |\alpha|\leq k} \ipl{D^{\alpha}u}{D^{\alpha}u}} = \hn{u}\\
    \norm{u}_{H^s}^2  \leq \hn{u}
\end{align}

\subsection{}

\begin{align}
    \hn{u+v}^2 &= \iph{u+v}{u+v}\\
               &= \iph{u}{u} + \iph{v}{v} + 2 \iph{u}{v}\\
    \hn{u+v}^2 &= \hn{u}^2 + \hn{v}^2 + 2 \sum_{|\alpha|\leq k}\ipl{D^{\alpha}u}{D^{\alpha}v}.
\end{align}

By Cauchy-Schwarz in the $L_2$ norm,

\begin{align}
    \hn{u+v}^2 &\leq \hn{u}^2 + \hn{v}^2 + 2 \sum_{|\alpha|\leq k}\norm{D^{\alpha}u}_{L_2(\Omega)} \norm{D^{\alpha}v}_{L_2(\Omega)}\\
    \hn{u+v}^2 &\leq \hn{u}^2 + \hn{v}^2 + 2 \sqrt{\left(\sum_{|\alpha|\leq k}\norm{D^{\alpha}u}_{L_2(\Omega)} \norm{D^{\alpha}v}_{L_2(\Omega)}\right)^2}.
\end{align}

Again by Cauchy-Schwarz,

\begin{align}
    \hn{u+v}^2 &\leq \hn{u}^2 + \hn{v}^2 + 2 \sqrt{\left(\sum_{|\alpha|\leq k} \norm{D^{\alpha}u}_{L_2(\Omega)}^2\right) \left(\sum_{|\alpha|\leq k} \norm{D^{\alpha}V}_{L_2(\Omega)}^2\right)}\\
    &\leq \hn{u}^2 + \hn{v}^2 + 2 \hn{u} \cdot \hn{v},
\end{align}

Taking the square root, this becomes 

\begin{equation}
    \hn{u+v} \leq \hn{u} + \hn{v}
\end{equation}


\end{document}